% TOP_PROBLEM: {"problemId": "problem.selberg-1-4-conjecture", "canonicalTitle": "Selberg 1/4 Conjecture", "releaseRank": 101, "primaryDomain": "number_theory_arithmetic_geometry", "primaryDomainLabel": "Number theory & arithmetic geometry", "displayStatus": "open", "releaseStatus": "open"}
% !TeX program = lualatex
% Definitions and sources retained from research_batch_101_103.tex; research is in GPT_6_Astra_Ultra.
\documentclass[11pt]{article}
\usepackage[margin=25mm]{geometry}
\usepackage{fontspec}
\setmainfont{Latin Modern Roman}
\usepackage{amsmath,amssymb,mathtools}
\usepackage{unicode-math}
\setmathfont{Latin Modern Math}
\usepackage{enumitem,needspace}
\usepackage{xurl,hyperref}
\hypersetup{colorlinks=true,urlcolor=blue}
\setcounter{secnumdepth}{1}
\setlength{\parindent}{0pt}
\setlength{\parskip}{5pt}
\setlength{\emergencystretch}{3em}
\setlist[itemize]{leftmargin=3em,itemsep=5pt,topsep=4pt,parsep=0pt}
\allowdisplaybreaks
\newcommand{\N}{\mathbb{N}}
\newcommand{\Z}{\mathbb{Z}}
\newcommand{\Q}{\mathbb{Q}}
\newcommand{\R}{\mathbb{R}}
\newcommand{\C}{\mathbb{C}}
\newcommand{\F}{\mathbb{F}}
\newcommand{\A}{\mathbb{A}}
\newcommand{\PP}{\mathbb P}
\newcommand{\E}{\mathbb{E}}
\newcommand{\eps}{\varepsilon}
\newcommand{\dd}{\,\mathrm d}
\newcommand{\sref}[2]{\hyperlink{source:#1:#2}{[S#2]}}
\newcommand{\eref}[2]{\hyperlink{extra:#1:#2}{[E#2]}}
% Turn the next switch off for a shorter reading copy. The quotations preserve
% every exactTarget field, including qualifications omitted from a short summary.
\newif\ifshowcatalogue
\showcataloguetrue
\newcommand{\cataloguescope}[1]{\ifshowcatalogue
 \paragraph{Exact scope in the supplied catalogue.}
 \begin{quote}\small #1\end{quote}\fi}
\begin{document}
\setcounter{section}{100}
% ================================================================
% problemId: problem.selberg-1-4-conjecture
% source releaseRank: 101; source releaseStatus: open
% exactTarget SHA-256: 369a061b0f8d65f7956e94416c425a610ae109ecd8cb0ba4744936d4982da2c1
\clearpage
\section{\texorpdfstring{Selberg 1/4 Conjecture}{Selberg 1/4 Conjecture}}\label{101}
\subsection{Definitions and mathematical statement}
The upper half-plane $\mathbb H=\{x+iy:y>0\}$ has hyperbolic metric $y^{-2}(dx^2+dy^2)$ and nonnegative Laplacian $\Delta=-y^2(\partial_x^2+\partial_y^2)$. A congruence subgroup contains a principal subgroup $\Gamma(N)$ for some $N$. On its finite-area quotient, the assertion excludes any nonconstant $L^2$ eigenfunction with
\[
 \Delta f=\lambda f,\qquad0<\lambda<\tfrac14.
\]
This formulation remains meaningful for noncompact quotients, where a continuous spectrum is also present; ``first nonzero eigenvalue'' is understood as the bottom of the nonconstant spectral range when necessary.
\par\smallskip\textit{Formulation and concept references:} \sref{101}{1}.
\cataloguescope{Prove that the first nonzero eigenvalue lambda\_1 of the Laplace-Beltrami operator on every congruence quotient of the upper half-plane is at least 1/4.}
\subsection{Short English statement}
Do all congruence hyperbolic surfaces have a spectral gap of at least one quarter above the constant functions?
\subsection{Sources}
\begin{itemize}\raggedright
\item[\textnormal{[S1]}]\hypertarget{source:101:1}{} Formal statement, Chapter 4, Selberg property.
\url{https://www.ma.huji.ac.il/~alexlub/BOOKS/On%20property/On%20property.pdf}.
{\small\textit{Catalogue formulation source.}}
\item[\textnormal{[S2]}]\hypertarget{source:101:2}{} Asymptotic independence for random permutations from surface groups.
\url{https://link.springer.com/article/10.1007/s10711-025-01003-8}.
{\small\textit{Catalogue research/status reference; not a proof endorsement.}}
\item[\textnormal{[S3]}]\hypertarget{source:101:3}{} Twist-minimal trace formulas and the Selberg eigenvalue conjecture.
\url{https://arxiv.org/abs/1803.06016}.
{\small\textit{Catalogue research/status reference; not a proof endorsement.}}
\item[\textnormal{[E1]}]\hypertarget{extra:101:1}{}
A. Pascadi, \emph{Large sieve inequalities for exceptional Maass forms and
the greatest prime factor of $n^2+1$}, Forum of Mathematics, Pi \textbf{14}
(2026), e8; arXiv:2404.04239v3, 15 January 2026. Introduction and Theorem C
(the Kim--Sarnak $7/64$ bound), and the exceptional-spectrum large-sieve
results. \url{https://arxiv.org/html/2404.04239v3}.
\url{https://doi.org/10.1017/fmp.2026.10025}.
{\small\textit{Consulted 25 September 2026; weighted estimates are not an
exclusion of exceptional eigenvalues.}}
\item[\textnormal{[E2]}]\hypertarget{extra:101:2}{}
G. Cherubini and M. S. Risager, \emph{On the variance of the error term in
the hyperbolic circle problem}, Revista Matem\'atica Iberoamericana
\textbf{34} (2018), no.~2, 655--685, DOI 10.4171/RMI/1000.
Section 2, transform convention and pre-trace formula (2.2); equation (1.2)
for the main term including small eigenvalues; Section 4 for smoothing and
absolute convergence.
\url{https://ems.press/content/serial-article-files/38698}.
{\small\textit{Consulted 25 September 2026.}}
% END RESEARCH ADDITION REFERENCES-101
\end{itemize}

\end{document}
